\documentclass[10pt,a4paper]{amsart}
\usepackage[margin=19mm]{geometry}
\usepackage{amsmath,amssymb,mathtools}
\usepackage[scr=rsfs,cal=euler]{mathalfa}
\usepackage{xfrac}
\usepackage[unicode,hidelinks,bookmarks=false]{hyperref}
\newcommand{\ie}{i.e.\ }
\newcommand{\cf}{cf.\ }
\newcommand{\resp}{respectively\ }
\newcommand{\Subsection}{\subsection}
\providecommand{\nobreakdash}{\nobreak\hbox{-}}
\setlength{\emergencystretch}{2em}
\multlinegap0pt
\usepackage{esint}
\usepackage{amssymb}
\usepackage{xcolor}
\usepackage{tikz}
\usepackage{subfig}
\usepackage{float}
\newtheorem{lemma}{Lemma}
\newtheorem{theorem}{Theorem}
\title{Regularity for Doubly Nonlinear Equations in the Mixed Regime}
\author{Simone Ciani, Eurica Henriques, Mariia Savchenko, Igor I. Skrypnik, Yevgeniia Yevgenieva}
\date{Source: arXiv:2602.09906v1 (CC BY 4.0)}
\begin{document}
\maketitle

\noindent\emph{Source and licence.} Regularity for Doubly Nonlinear Equations in the Mixed Regime. Version arXiv:2602.09906v1 (CC BY 4.0), distributed by arXiv under the Creative Commons Attribution 4.0 International licence, http://creativecommons.org/licenses/by/4.0/, as recorded in the arXiv licence metadata for this exact version and shown on the abstract page https://arxiv.org/abs/2602.09906v1. Full licence notice: \texttt{licenses/arxiv-2602.09906-CC-BY-4.0.txt}.\par\smallskip

\noindent\textit{Editorial scope.} Counted unit: the article's theorem lem3.1 (source0.tex lines 820--827). The article's complete original proof of it is delivered with it (source0.tex lines 828--933, one \texttt{proof} environment of 106 lines). No mathematics is written, repaired or re-derived: the statement and the proof are the article's own lines, in the article's own notation, and no retained line is substituted or re-flowed.  Retained uncounted as the source-local context the proof needs: lines 179--181; lines 186--191; lines 201--204; lines 744--754.  Nothing was omitted from any retained span, so the package carries no omission record.  No retained line contains a citation, so the wrapper carries no bibliography and no bibliography entry.  No retained line contains a figure environment, an image-inclusion command or drawing code, so no image is delivered and the package declares no asset dependency.  Editorial formatting only, in addition: an AMS \texttt{amsart} preamble that restates the article's own declarations and macros used by the retained lines, namely the environments \texttt{lemma}, \texttt{theorem} on separate counters, together with the standard packages those lines need.  The retained environments print in this document as Lemma 1, Theorem 1 in order of appearance, whereas the article prints the same items with the numbering of its own full document.  The complete native source of the article as distributed in the arXiv e-print for this exact version is bundled unchanged as \texttt{sources/194354/source0.tex}.
\begin{equation}\label{eq1.1}
u_{t}-\textrm{div}\mathbf{A}(x, t, u, D u)=0, \quad \quad  (x, t)\in \Omega_{T}.
\end{equation}

\begin{equation}\label{eq1.2}
\begin{cases}
\mathbf{A}(x, t, u, D u) D u^{m^-} \geqslant  K_{1}\,u^{(m-m^-)(p-1)}|D u^{m^-}|^{p},\\
|\mathbf{A}(x, t, u,  D u)| \leqslant  K_{2}\,u^{(m-m^-)(p-1)}|D u^{m^-}|^{p-1},
\end{cases}
\end{equation}

\begin{equation}\label{eq1.3}
\int\limits_{E}u \zeta\, dx \bigg|^{t_{2}}_{t_{1}} + \int\limits^{t_{2}}\limits_{t_{1}}\int\limits_{E}
\{-u\zeta_{\tau}+ \mathbf{A}(x, \tau, u, D u) \,D \zeta\}\, dx d\tau \leqslant (\geqslant)\, 0,
\end{equation}

\begin{lemma}\label{lem2.3}
Let $u$ be a nonnegative, local weak supersolution of \eqref{eq1.1}-\eqref{eq1.2} in $\Omega_T$. Fix parameters $0<\alpha<1$ and $1<\beta<1+\frac{1}{m^{-}}$. Then there exists a constant $\gamma>0$, depending only on the data and on $\beta$, such that for every $\varepsilon>0$, every cylinder $B_r(y)\times(s,t)\subset\Omega_T$, and every piecewise smooth cutoff function $\zeta$ vanishing on $\partial B_r(y)$ and satisfying $0\leqslant\zeta\leqslant 1$, the following inequality holds:
\begin{equation}\label{eq2.2}
\begin{aligned}
&\int\limits_s^t\int\limits_{B_r(y)}\Big(\frac{\tau-s}{t-s}\Big)^\alpha
\frac{u^{(m-m^-)(p-1)}}{(u^{m^-}+\varepsilon^{m^-})^\beta}|D u^{m^-}|^{p}\,\zeta^{p}\,dx d\tau
\\&\leqslant \gamma \int\limits_{B_r(y)\times\{t\}}(u+\varepsilon)^{1-m^-(\beta-1)}\,dx+
\gamma \int\limits_s^t\int\limits_{B_r(y)}(u+\varepsilon)^{\lambda-m^-(\beta-1)}\,|D \zeta|^{p}\,dx d\tau.
\end{aligned}
\end{equation}
\end{lemma}

\begin{theorem}\label{lem3.1}
Let $u$ be a nonnegative, local weak solution to \eqref{eq1.1}-\eqref{eq1.2}, and assume that $\lambda<1$. Then there exists a constant $\gamma>0$, depending only on the data, such that for any cylinder $Q_{2r,t-s}(y,t)\subset\Omega_T$, the following estimate holds:
\begin{equation}\label{eq3.1}
\sup\limits_{s\leqslant \tau\leqslant t}\fint\limits_{B_{r}(y)}u(x, \tau)\, dx \leqslant \gamma
\inf\limits_{s\leqslant \tau \leqslant t}\fint\limits_{B_{2 r}(y)} u(x, \tau)\,dx +\gamma
\Big(\frac{t-s}{r^{p}}\Big)^{\frac{1}{1-\lambda}}.
\end{equation}
\end{theorem}

\begin{proof}
Without loss of generality, we set $s=0$.
For $j=0,1,2,\dots$, define
$$
r_j:= \sum\limits^{j}\limits_{l=1} \frac{r}{2^{l}},
\qquad
\bar{r}_j := \frac{r_j+ r_{j+1}}{2},
$$
and set
$$
B_j:=B_{r_j}(y),
\qquad
\bar B_j:=B_{\bar{r}_j}(y),\qquad B_j\subset\bar B_j\subset B_{j+1}.
$$
Let $\zeta_j\in C_0^1(\bar B_j)$ be a cutoff function satisfying
$$
\zeta_j\equiv1 \text{ in } B_j,
\qquad
0\leqslant\zeta_j\leqslant 1,
\qquad
|D_i\zeta_j|\leqslant \gamma\,2^{j}r^{-1},
\quad i=1,\dots,N.
$$
Fix $0<t_1<t_2<t$.
Testing \eqref{eq1.3} with $\zeta_j(x)$ and integrating over
$(t_1,t_2)\times\bar B_j$, we obtain
\begin{equation}\label{pre-recurs_ineq}
\fint\limits_{\bar B_j}u(x,t_2)\zeta_j\,dx
\leqslant
\fint\limits_{\bar B_j}u(x,t_1)\zeta_j\,dx
+
\frac{\gamma\,2^j}{r}
\int\limits_{t_1}^{t_2}\fint\limits_{\bar B_j}
u^{(m-m^-)(p-1)}
|D u^{m^-}|^{p-1}\,dx\,d\tau .
\end{equation}
Choose $t_1$ such that
$$
\fint\limits_{B_{2r}(y)} u(x, t_{1})\,dx= \inf\limits_{0<\tau<t} \fint\limits_{B_{2r}(y)\times\{\tau\}} u\,\,\,dx,
$$
Since $\bar B_j\subset B_{2 r}$, we obtain
$$
\fint\limits_{\bar B_j}u(x,t_1)\,\zeta_j\,dx
\leqslant \gamma
\inf\limits_{0<\tau<t} \fint\limits_{B_{2r}(y)\times\{\tau\}} u\,\,\,dx.
$$
Define $J_j:=\sup\limits\limits_{0<\tau<t}\,\,\fint\limits_{B_j}u(x,\tau)\,dx$, we arrive at
\begin{equation}\label{recurs_1}
J_j
\leqslant \gamma
\inf\limits_{0<\tau<t}\fint\limits_{B_{2r}(y)\times\{\tau\}}u\,\,\,dx
+
\frac{\gamma\,2^j}{r}
\int\limits_0^t\fint\limits_{\bar B_j}
u^{(m-m^-)(p-1)}
|D u^{m^-}|^{p-1}\,dx\,d\tau .
\end{equation}
In order to estimate the second term in the previous inequality, we fix  $\varepsilon =  \left(\frac{t}{r^{p}}\right)^{\frac{1}{1-\lambda}},$
and let $\alpha>0$ and $\beta>1$ be constants to be chosen later.  
Applying H\"older's inequality, we obtain
\begin{align}\label{eq5.4}
&\frac{\gamma\,2^j}{r}\int\limits_0^{t}\fint\limits_{\bar B_j}u^{(m-m^-)(p-1)}|D u^{m^-}|^{p-1}\,dx d\tau
\leqslant  \frac{\gamma\,2^j}{r} \Big(\int\limits_0^{t}\fint\limits_{ B_{j+1}}\Big(\frac{\tau}{t}\Big)^{\alpha}\frac{u^{(m-m^-)(p-1)}}{(u^{m^{-}}+\varepsilon^{m^{-}})^{\beta}}|D u^{m^-}|^{p}\,dx d\tau\Big)^{\frac{p-1}{p}}\notag\\
&\quad\times\Big(\int\limits_0^{t}\fint\limits_{B_{j+1}}
\Big(\frac{t}{\tau}\Big)^{\alpha(p-1)} u^{(m-m^{-})(p-1)}(u^{m^{-}}+\varepsilon^{m^{-}})^{\beta(p-1)}\,dx\,d\tau\Big)^{\frac{1}{p}}= \frac{\gamma 2^j}{r}  A_1^{\frac{p-1}{p}} A_2^{\frac{1}{p}}.
\end{align}
In order to estimate the right-hand side of \eqref{eq5.4}, we start by applying the elementary inequality for positive numbers: $u^{m^-(\frac{m}{m^-}-1)(p-1)} 
\leqslant(u^{m^-}+\varepsilon^{m^-})^{(\frac{m}{m^-}-1)(p-1)}, \varepsilon>0, $ and then take the supremum of the spatial average over the time interval $[0,t]$, to get
\begin{align*}
A_2&\leqslant\int\limits_0^{t}\fint\limits_{B_{j+1}}\Big(\frac{t}{\tau}\Big)^{\alpha(p-1)} (u^{m^{-}}+\varepsilon^{m^{-}})^{(\frac{m}{m^{-}}+\beta -1)(p-1)} dx d\tau
\\&\leqslant \sup\limits_{0<\tau<t}\fint\limits_{B_{j+1}}(u+\varepsilon)^{(m+m^{-}(\beta -1))(p-1)} dx \ \int\limits_0^t \Big(\frac{t}{\tau}\Big)^{\alpha(p-1)} d\tau.
\end{align*}
In what follows, we fix the values $\alpha$ and $\beta$. By choosing $\alpha$ such that $\alpha(p-1)<1$, the time integral on the right-hand side of the previous inequality is finite and can be estimated as
$\int_0^t \Big(\frac{t}{\tau}\Big)^{\alpha(p-1)}\, d\tau 
\leqslant \gamma\, t,$
where $\gamma$ is a positive constant depending only on $\alpha$ and $p$. Furthermore by choosing $\beta$ so that $(m+m^-(\beta-1))(p-1)<1$,  we can apply H\"older's inequality to the spatial integral. Hence,
\begin{equation}\label{eq5.5}
A_2\leqslant \gamma \,t \,\Big(\sup\limits_{0<\tau<t}\fint\limits_{ B_{j+1}\times \{\tau\}}(u+\varepsilon) \,dx\Big)^{(m+m^{-}(\beta -1))(p-1)}.
\end{equation}
To estimate $A_1$ of \eqref{eq5.4}, we apply Lemma~\ref{lem2.3}. 
Choosing $\beta$ such that $m^{-}(\beta-1)<1$,  using the simple estimate   $(u+\varepsilon)^{\lambda-1}\leqslant \varepsilon^{\lambda-1}$. Recalling the choice of $\varepsilon$
 and applying H\"older's inequality, we obtain
\begin{equation}\label{eq5.6}
\begin{aligned}
A_1&\leqslant \gamma\fint\limits_{B_{j+1}\times\{t\}}(u+\varepsilon)^{1-m^-(\beta-1)} dx+\gamma \int\limits_0^t\fint\limits_{B_{j+1}}
(u+\varepsilon)^{\lambda-m^-(\beta-1)}|D \zeta_{j}|^{p}\,dxd\tau\\
&\leqslant \gamma 2^{jp}\bigg[1+ \frac{t}{r^{p}\varepsilon^{1-\lambda}}\bigg]\sup\limits_{0<\tau<t}\fint\limits_{B_{j+1}\times\{\tau\}}
(u+\varepsilon)^{1-m^-(\beta-1)} dx\\&\leqslant \gamma 2^{jp}
\Big(\sup\limits_{0<\tau<t}\fint\limits_{B_{j+1}\times\{\tau\}}(u+\varepsilon)\,\, dx\Big)^{1-m^-(\beta-1)}.
\end{aligned}
\end{equation}
Combining \eqref{eq5.4}--\eqref{eq5.6} and applying Young's inequality with $\delta \in (0,1)$ and exponents  $\frac{p}{p-(1-\lambda)}$, $\frac{p}{1-\lambda}$, we derive 
\begin{align*}
&\frac{\gamma\,2^j}{r}\int\limits_0^{t}\fint\limits_{\bar B_{j}}u^{(m-1)(p-1)}|D u|^{p-1}\,dx d\tau
\leqslant \gamma 2^{j \gamma} \Big(\frac{t}{r^{p}}\Big)^{\frac{1}{p}}\Big(\sup\limits_{0<\tau<t}\fint\limits_{B_{j+1}\times\{\tau\}}(u+\varepsilon) \,\,dx\Big)^{\frac{(p-1)(m+1)}{p}}\\
&\leqslant
\delta \sup\limits_{0<\tau<t}\fint\limits_{B_{j+1}\times\{\tau\}}(u+\varepsilon)\,\, dx+\frac{\gamma 2^{j\gamma}}{\delta^{\gamma}}\Big(\frac{t}{r^{p}}\Big)^{\frac{1}{1-\lambda}}\leqslant \delta \sup\limits_{0<\tau<t}\fint\limits_{B_{j+1}\times\{\tau\}}(u+\varepsilon)\,\, dx+\frac{\gamma 2^{j\gamma}}{\delta^{\gamma} } \varepsilon,
\end{align*}
where $\gamma$ also depends on $\alpha$ and $\beta$, which are chosen as $\alpha=\min\Big(\frac{1}{2}, \frac{1}{2(p-1)}\Big)$ and $\beta=1+\frac{1}{2m^{-}} \min\big(1, \frac{1-\lambda}{p-1}\big)$.
Combining this estimate with \eqref{recurs_1}, we obtain
\begin{equation*}
J_{j} \leqslant \delta J_{j+1} + \gamma \delta^{-\gamma} 2^{j \gamma}\big(\inf\limits_{0< \tau <t} \fint\limits_{B_{2r}(y)\times\{\tau\}} u\ dx
+ \varepsilon\big), \quad j=0, 1, 2,\ldots
\end{equation*}
Iterating this inequality and choosing $\delta$ sufficiently small, we arrive at \eqref{eq3.1}.
\end{proof}

\end{document}
